\documentclass[UTF8,zihao=-4]{ctexart}
\usepackage[a4paper,margin=2.2cm]{geometry}
\usepackage{fontspec}
\usepackage{tabularx}
\usepackage{array}
\usepackage{ragged2e}
\usepackage{fancyhdr}
\usepackage{hyperref}
\setCJKmainfont{Noto Serif SC}
\setCJKsansfont{Noto Sans SC}
\setmainfont{Noto Serif SC}
\setlength{\parindent}{0pt}
\setlength{\parskip}{0.45em}
\pagestyle{fancy}
\fancyhf{}
\fancyhead[L]{商务智能}
\fancyfoot[C]{\thepage}
\hypersetup{colorlinks=true,linkcolor=black,urlcolor=blue}
\begin{document}
\title{7、数据仓库的特征}
\author{}
\date{}
\maketitle

数据仓库的四大特征就是：面向主题、集成、非易失、时变。\par

1.1 面向主题\par

数据仓库不是按照具体业务应用或部门系统来组织数据，而是按照管理决策中的分析主题来组织数据。\par

例如主题可以是：\par

商品\par
供应商\par
顾客\par
CRM 中的客户关系管理\par
ERP 中的销售管理、产品质量控制、库存管理等\par


1.2 集成\par

集成是指数据仓库要把来自不同部门、不同系统、不同平台、不同格式的数据统一起来。\par

操作型系统中的数据可能存在：\par

同名异义\par
同义异名\par
单位不一致\par
编码不一致\par
数据类型不一致\par
数据分散在不同系统中\par

数据仓库通过 ETL 把它们清洗、转换、统一后形成一致的数据环境。\par

1.3 非易失 / 稳定\par

数据仓库中的数据主要用于分析处理，通常不是实时修改，而是以查询、分析为主。PPT 中也把分析型数据的操作类型概括为“读，稳定的”。\par


1.4 时变\par

数据仓库保存历史数据，具有时间跨度。操作型数据库主要保存当前数据，而数据仓库关注历史变化，用于研究趋势。PPT 中指出，分析处理更看重 5 到 10 年的历史数据，而事务处理一般只需要当前或短期数据。\par

回答：\par
1.面向主题,指数据仓库中的数据按照较高层次的分析对象进行组织，而不是按照具体应用系统组织。主题是企业信息系统中数据的综合、归类和分析利用对象。\par
2.集成，集成是指数据仓库要把来自不同部门、不同系统、不同平台、不同格式的数据统一起来。数据仓库通过 ETL 把它们清洗、转换、统一后形成一致的数据环境。\par
3.非易失，非易失性是指数据进入数据仓库后一般不再像操作型数据库那样频繁更新、删除或修改，而是相对稳定地保存，用于查询和分析。\par
4.时变，时变性是指数据仓库中的数据反映一段时间内的数据变化，通常带有时间属性，可以支持趋势分析和历史分析。\par

\end{document}
